% ==============================================================================
% HALAMAN SAMPUL (COVER)
% Sesuai Lampiran 1b Pedoman Penulisan Skripsi FST UNIBA Madura 2026
% ==============================================================================

\begin{titlepage}
\begin{center}
    \singlespacing
    
    % Judul Skripsi (Kapital, Tebal, 12pt proporsional)
    {\fontsize{12pt}{16pt}\selectfont \bfseries \skripsiJudul}
    
    \vspace{1.5cm}
    
    % Tulisan Tipe Dokumen (PROPOSAL SKRIPSI / SKRIPSI) & Bidang Minat
    {\fontsize{12pt}{16pt}\selectfont \bfseries \MakeUppercase{\skripsiTipeDokumen}}\\[0.25cm]
    {\fontsize{12pt}{16pt}\selectfont \skripsiBidangMinat}
    
    \vspace{0.9cm}
    
    % Kalimat Persyaratan (12pt reguler)
    {\fontsize{12pt}{16pt}\selectfont
    \skripsiTujuanDokumen
    }
    
    \vspace{1.2cm}
    
    % Logo Universitas Bahaudin Mudhary Madura
    \includegraphics[width=4.0cm]{figures/logo-uniba.png}
    
    \vspace{1.2cm}
    
    % Identitas Mahasiswa (Nama tebal, NIM reguler)
    {\fontsize{12pt}{16pt}\selectfont \bfseries \MakeUppercase{\skripsiPenulis}}\\
    \vspace{0.2cm}
    {\fontsize{12pt}{16pt}\selectfont NIM. \skripsiNIM}
    
    \vfill
    
    % Identitas Institusi
    {\fontsize{12pt}{16pt}\selectfont \bfseries
    PROGRAM STUDI \MakeUppercase{\skripsiProdi}\\
    \MakeUppercase{\skripsiFakultas}\\
    \MakeUppercase{\skripsiUniversitas}\\
    \MakeUppercase{\skripsiKota}\\
    \skripsiTahun
    }

\end{center}
\end{titlepage}
